\documentclass[11pt, a4paper]{article}

%----------
% To automatically save paper and trees (at least, for preprints):
\usepackage{fullpage} % or \usepackage[headings]{fullpage}
%----------

%----------
% To load external graphics:
\usepackage{graphicx}
%----------

%----------
% AMS packages, very recommended for maths:
\usepackage{amsmath, amsthm, amsfonts, amssymb}
% \usepackage{mathtools}
\allowdisplaybreaks[3]
%----------

%----------
\newcommand{\N}{\mathbb{N}}
\newcommand{\Z}{\mathbb{Z}}
\newcommand{\C}{\mathbb{C}}
%----------

%----------
\renewcommand{\Re}{\operatorname{Re}}
\renewcommand{\Im}{\operatorname{Im}}
\newcommand{\calF}{\mathcal{F}}
\newcommand{\p}{\mathcal{P}} % \P ya existe
\newcommand{\manF}{\Lambda_{\mathcal{F}}} 
%\newcommand{\uF}{u_\mathcal{F}} 
%\newcommand{\muF}{\mu_\mathcal{F}} 
\newcommand{\Fp}{\mathcal{F}_{\mathcal{P}}} 
\newcommand{\IF}{\operatorname{IF}} 
%----------

%----------
\theoremstyle{plain}
\newtheorem{theorem}{Theorem}
\newtheorem{proposition}[theorem]{Proposition}
\newtheorem{lemma}[theorem]{Lemma}
%\newtheorem{corollary}[theorem]{Corollary}
\theoremstyle{definition}
\newtheorem{remark}[theorem]{Remark}
%----------

% %----------
% %%% Para BORRADORES:
% \usepackage{color}
% \usepackage[svgnames,x11names,dvipsnames]{xcolor}
% \newcommand{\rojo}[1]{\textcolor{red}{#1}}
% \newcommand{\azul}[1]{\textcolor{blue}{#1}}
% \newcommand{\verde}[1]{\textcolor{green}{#1}}
% \usepackage[notcite,notref]{showkeys}
% %----------

%----------
% To automatically include internal hyperlinks 
% and internet links with \url{http://....} 
\usepackage{hyperref}
\usepackage{orcidlink}
%----------

%----------
%%% Para BORRADORES:
%\usepackage[notcite,notref]{showkeys}
%\usepackage{refcheck}
%\usepackage{backref}
%----------

%% Gaspar Mora:
%% Facultad de Ciencias Exactas y Naturales, Universidad Nacional de Asunci\'on, 
%% Campus Universitario de San Lorenzo, San Lorenzo, Paraguay
%% gaspar.mora@facen.una.py
%% https://orcid.org/0000-0003-0883-220X
%% https://www.facen.una.py

%% Luis M. Navas:
%% navas@usal.es
%% Luis Manuel Navas Vicente
%% Departamento de Matem\'aticas and Instituto Universitario de F\'isica Fundamental y Matem\'aticas (IUFFyM),
%% Universidad de Salamanca, 
%%%% Plaza de la Merced, 1--4,
%% 37008 Salamanca, Spain
%% https://orcid.org/0000-0002-5742-8679

%% Juan L. Varona:
%% Departamento de Matem\'aticas y Computaci\'on, 
%% Universidad de La Rioja, 26006 Logro\~no, Spain
%% jvarona@unirioja.es
%% https://orcid.org/0000-0002-2023-9946

%% Luis y Juan Luis:
%% Partially supported by grant PID2021-124332NB-C22 % antiguo
%% (Ministerio de Ciencia e Innovaci\'on-Agencia Estatal de Investigaci\'on, Spain).
%% PID2024-155593NB-C22 % nuevo
%% (Ministerio de Ciencia, Innovaci\'on y Universidades-Agencia Estatal de Investigaci\'on, Spain).

%% Luis: 
%% Partially supported by grant PID2023-150787NB-I00

%----------
\title{On the von Mangoldt-type function\\of the Fibonacci zeta function\thanks{%
This research was partially supported by grants PID2023-150787NB-I00 and PID2024-155593NB-C22
(Minis\-te\-rio de Cien\-cia, Inno\-va\-ci\'on y Uni\-ver\-si\-da\-des, Spain).}
\textsuperscript{,}\thanks{\textbf{This paper has been published in:} \textit{Rev.\ R.\ Acad.\ Cienc.\ Exactas\ F\'is.\ Nat.\ Ser.\ A Mat.\ RACSAM} \textbf{119} (2025), no.~4, Paper No.~126, 19~pp., \url{https://doi.org/10.1007/s13398-025-01792-3}}}
%----------

%----------
%\author{Gaspar Mora, Luis M. Navas and Juan L. Varona}
\author{Gaspar Mora\textsuperscript{1}\,\orcidlink{0000-0003-0883-220X}, 
Luis M. Navas\textsuperscript{2}\,\orcidlink{0000-0002-5742-8679}, 
and Juan L. Varona\textsuperscript{3}\,\orcidlink{0000-0002-2023-9946}\\[6pt]
  \small 
  \textsuperscript{1}Facultad de Ciencias Exactas y Naturales, Universidad Nacional de Asunci\'on,\\[-2pt]
  \small Campus Universitario de San Lorenzo, San Lorenzo, Paraguay\\[-1pt]
  \small {gaspar.mora@facen.una.py}\\[3pt]
  \small 
  \textsuperscript{2}Departamento de Matem\'aticas and 
  Instituto Universitario de \\[-2pt]
  \small F\'isica Fundamental y Matem\'aticas (IUFFyM), Universidad de Salamanca, Spain\\[-1pt]
  \small {navas@usal.es}\\[3pt]
  \small 
  \textsuperscript{3}Departamento de Matem\'aticas y
  Computaci\'on, Universidad de La Rioja,\\[-2pt]
  \small Logro\~no, Spain\\[-1pt]
  \small {jvarona@unirioja.es}
}
%----------

%----------
\date{}     % Activate to display a given date or no date
%----------

%----------
\begin{document}
%----------

%----------
\maketitle
%----------

\begin{abstract}
The Dirichlet series associated to the Fibonacci sequence $\{F_{n}\}$, 
\begin{equation*}
  \sum_{n=1}^{\infty} F_{n}^{-s},
  %\quad s\in \C, \quad \Re s>0,
\end{equation*}
converges for $s\in \mathbb{C}$ with $\Re s > 0$. The analytic function $\varphi(s)$ it defines on the right half-plane is known as the Fibonacci zeta function. Here we consider its logarithmic derivative $\varphi'(s)/\varphi(s)$, which formally corresponds to the Dirichlet series 
\begin{equation*}
  % \frac{\varphi'(s)}{\varphi(s)} \leftrightsquigarrow
  -\sum_{l=1}^{\infty} \Lambda_{\mathcal{F}}(l)l^{-s},
\end{equation*}
where the arithmetical function $\Lambda_{\mathcal{F}}(l)$ can be considered analogous to the classical von Mangoldt function $\Lambda(s)$, which is defined by $\zeta'(s)/\zeta(s) = -\sum_{n=1}^{\infty} \Lambda(n) n^{-s}$ where $\zeta(s)$ is the Riemann zeta function. This paper studies some properties of the function $\Lambda_{\mathcal{F}}(l)$ along with the domain of convergence of this Dirichlet series.

\medskip

\textbf{MSC2020:} Primary 11M41; Secondary 11A25, 11B39, 30B50.
%11A25 - Arithmetic functions; related numbers; inversion formulas
%11B39 - Fibonacci and Lucas numbers and polynomials and generalizations
%11M41 - Other Dirichlet series and zeta functions
%30B50 - Dirichlet series, exponential series and other series in one complex variable

\smallskip

\textbf{Keywords:} Fibonacci numbers, Fibonacci zeta function, Dirichlet series, von Mangoldt-Fibonacci function.
\end{abstract}

%----------
\section{Introduction and main results}
%----------

Let us denote the golden ratio by $\phi = \frac{1+\sqrt{5}}{2}$. It is well known that the $n$\textsuperscript{th} Fibonacci number, $F_{n}$, defined by the recurrence relation
\[
  F_{n+2} := F_{n+1}+F_{n}, \quad n\geq 1,\qquad F_1=1,\quad F_2=1, 
\]
can be expressed in terms of $\phi$ and $\psi := 1-\phi = -1/\phi = \frac{1-\sqrt{5}}{2}$ by means of Binet's formula 
%% Table[FullSimplify[(GoldenRatio^n - (-GoldenRatio)^(-n))/Sqrt[5]], {n, 1, 18}]
%% {1, 1, 2, 3, 5, 8, 13, 21, 34, 55, 89, 144, 233, 377, 610, 987, 1597, 2584}
%%
\begin{equation*}
%\label{eq:Binet}
F_{n} = \frac{\phi^{n} - \psi^n}{\phi - \psi}
= \frac{\phi^{n} - (-\phi)^{-n}}{\sqrt{5}},
\qquad n\geq 1.
\end{equation*}
The Fibonacci zeta function is defined by the Dirichlet series
\begin{equation}
\label{eq:phiFib}
\varphi(s) := \sum_{n=1}^{\infty} F_{n}^{-s},\quad s\in \C,
\quad \Re s>0
\end{equation}
(the convergence for $\Re(s)>0$ is clear from the asymptotic relation $F_{n} \approx \phi^{n}/\sqrt{5}$ when $n\to\infty$). Around the turn of the millennium Egami \cite{Ega} and Navas \cite{Na} independently obtained the meromorphic continuation of \eqref{eq:phiFib} to the complex plane, via the series
\[
\varphi(s) = 5^{s/2} \sum_{k=0}^{\infty}
\binom{-s}{k} \frac{1}{\phi^{s+2k}+(-1)^{k+1}}, 
\]
which exhibits simple poles at the points 
\[
s = -2k+\frac{\pi i(2m+k)}{\log \phi},
\quad k, m\in \Z,
\text{ with }k\geq 0.
\]
Later, Ram Murty \cite{Mur} easily derived the analytic continuation of \eqref{eq:phiFib} by using a technique that he had previously developed jointly with
Sinha~\cite{MurSin}. See also the recent paper \cite{SeNa} for the analytic continuation of general Dirichlet $\sum_{n=1}^{\infty} a_{n}^{-s}$ where $a_{n}$ satisfies a linear recurrence of arbitrary degree with integer coefficients.

As for the zeros of $\varphi(s)$, in \cite[Corollary~7]{Mo} it was proved that $\varphi(s) \neq 0$ for all $s$ in the half-plane 
\[
  H_{\eta} = \{ s\in \C :\Re s>\eta \}, 
\]
where $\eta $ is the unique positive real number satisfying $\varphi(\eta )=4$, and whose approximate value is 
\[
\eta = 0.7570549496906548985355124\dots
\]
(see \cite[Theorem~13]{Mo}). Thus $\frac{\varphi'(s)}{\varphi(s)}$ is analytic on $H_{\eta}$.
We are interested in the formal Dirichlet series associated to this function, which we will denote by
\begin{equation}
\label{eq:ManFib}
  -\sum_{l=1}^{\infty} \manF(l)l^{-s},
\end{equation}
where the arithmetical function
\[
  \manF : \N \rightarrow \C
\]
is, in this context, analogous to the classical von Mangoldt function, defined by 
\[
  \Lambda (l) := 
  \begin{cases}
    \log(p), & \text{ if } l = p^{k}, \ p\text{ prime, }k\geq 1, 
         \\ 
        0, & \text{ otherwise,}
  \end{cases}
\]
and which is determined by the formal relation (see, for instance, \cite[p.~3]{Ivic})
\[
  \frac{\zeta'(s)}{\zeta(s)} = -\sum_{l=1}^{\infty} \Lambda (l)l^{-s},
  \quad \Re s>1. 
\]
We will study the values $\manF(l)$ and the convergence of~\eqref{eq:ManFib}. 

Since $\varphi'(s) = - \sum_{n=1}^\infty \log(F_{n}) F_{n}^{-s}$, the relation
\[
  \varphi'(s) = \varphi(s) \cdot \frac{\varphi'(s)}{\varphi(s)}  
\]
means that
\begin{equation}
\label{eq:producto}
\sum_{n=1}^\infty \log(F_{n}) F_{n}^{-s} 
= \left( \sum_{n=1}^\infty F_{n}^{-s} \right)
\left( \sum_{l=1}^\infty \manF(l)l^{-s} \right).
\end{equation}
Recall that the formal product of two Dirichlet series $f(s) = \sum_{n=1}^\infty a(n)n^{-s}$ and $g(s) = \sum_{n=1}^\infty b(n)n^{-s}$ corresponds to the Dirichlet series $h(s) = f(s)g(s) = \sum_{n=1}^\infty c(n)n^{-s}$ of the Dirichlet convolution  $c(n) = (a*b)(n) = \sum_{rs=n} a(r)b(s)$ of their coefficients (and moreover, if $f(s)$ is absolutely convergent for $\Re(s)>\sigma_f$ and $g(s)$ for $\Re(s)>\sigma_g$, then $h(s)$ is absolutely convergent for $\Re(s) > \max\{\sigma_f, \sigma_g\}$); see, for instance, \cite[Chapter~11]{Ap}. 
In the case of \eqref{eq:producto}, this allows us to recursively compute the coefficients $\manF(l)$ using~\eqref{eq:producto}. In particular, the formal relation
\begin{equation}
\label{eq:ManDir}
\frac{\varphi'(s)}{\varphi(s)} = -\sum_{l=1}^{\infty} \manF(l)l^{-s}
\end{equation}
will be valid analytically if the right-hand-side series converges in some half-plane. 
The main goal of this paper is to determine $\sigma_0>0$ such that this series converges absolutely for $\Re(s)>\sigma_0$. Along the way, we will prove some interesting properties of the function $\manF(l)$. As far as we know, these functions have not been previously considered in the mathematical literature; see \cite{GoSh}, which collects many interesting Dirichlet series, as well as \cite[Chapters 16 and~17]{HaWr} and  \cite[Chapters 1 and~2]{Kno}. 

Since we will be considering only the multiplicative structure of the Fibonacci numbers, we only need the those that are greater than $1$, that is,
\[
\calF = \{ F_{n} : n\ge3 \} = \{ 2, 3, 5, 8, 13, 21, 34, 55, 89, 144, 233, 377, 610,\dots \}.
\]
We denote the multiplicative semigroup generated by $\calF$ by
\begin{equation}
\label{eq:Fpbasic}
\Fp := \{\text{finite products of powers of numbers in $\calF$}\};
\end{equation}
% sub\'indice ``$\p$'' alude a ``producto''
and we assume $1 \in \Fp$, by thinking of it as the empty product. In particular, $\calF \subseteq \Fp$.

To generate all the numbers in \eqref{eq:Fpbasic}, it is not necessary to use $F_6=8$ and $F_{12}=144$, because they can be written as product of other Fibonacci numbers, namely $F_6 = 8 = F_3^3$ and $F_{12} = 144 = 2^4 \cdot 3^2 = F_3^4 F_4^2$ (besides $F_1=F_2=1$, these are also the unique Fibonacci numbers that are perfect powers, a non-trivial result proved in~\cite{BMS}). 

Furthermore, Carmichael's Theorem (see \cite[Theorem~23, p.~61]{Car} or, for a simpler proof, \cite[Theorem 3]{Ya}) states that, if $n \ne 1,2,6,12$, the Fibonacci number $F_{n}$ contains at least one primitive divisor, defined as a prime divisor $p$ such that $p$ does not divide any $F_m$ for $0<m<n$. Hence, $F_{6}$ and $F_{12}$ are the only Fibonacci numbers expressible as a power or as a product of powers of smaller Fibonacci numbers, so that
\begin{equation}
\label{eq:Fp}
\Fp = \left\{
l = F_{n_1}^{m_1}F_{n_2}^{m_2} \cdots F_{n_{k_l}}^{m_{k_l}} :
m_{j}>0, \ 3 \le n_1<n_2<\cdots<n_{k_l}, \ n_{j}\neq 6, 12 \right\}, 
\end{equation}
where $k_{l}$ is the number of different Fibonacci numbers in $\calF \setminus \{8,144\}$ that appear in the decomposition of $l$; again, we assume $1 \in \Fp$ is the product of zero factors. 

By imitating the standard proof of unique factorization of natural numbers, it is a straightforward exercise to check that Carmichael's Theorem implies unique factorization in $\Fp$. When we refer to ``the'' factorization of a number in $\Fp$, this is the one we mean.

We now present our list of results concerning the function $\manF(l)$ and the convergence of the Dirichlet series \eqref{eq:ManDir}. The corresponding proofs will be given in the following sections.

The first result concerns the support of $\manF(l)$:

\begin{proposition}
\label{prop:Man(l)=0}
Let $\manF(l)$ be the arithmetical function defined in~\eqref{eq:producto}. Then, $\manF(1)=0$ and $\manF(l) = 0$ when $l \notin \Fp$.
\end{proposition}

As a consequence of this result, the Dirichlet series \eqref{eq:ManDir} can be written as
\begin{equation}
\label{eq:ManFp}
\frac{\varphi'(s)}{\varphi(s)}
= -\sum_{l\in \Fp} \manF(l)l^{-s}.
\end{equation}

For $l$ equal to a product of powers of (ordered) Fibonacci numbers excluding $F_3=2$, or including $F_3$ but raised to a power less than $3$, the expression for $\manF(l)$ is the following:

\begin{proposition}
\label{prop:Man(l)prod}
Let $\manF(l)$ be the arithmetical function defined in~\eqref{eq:producto}. 
Then, if 
\[
l = F_{n_1}^{m_1} F_{n_2}^{m_2} \cdots F_{n_{k_l}}^{m_{k_l}} \in \Fp
\] 
(with a factorization as in \eqref{eq:Fp}, which is in fact unique) 
with $F_{n_1}\neq 2$ or $F_{n_1}=2$ and $m_1<3$, one has 
\begin{equation}
\label{eq:Man(l)prod}
\manF(l) =
(-1)^{r_l-1}\frac{1}{2^{r_{l}}}\frac{(r_l-1)!}{m_1!\, m_2! \cdots m_{k_l}!}\log(l), 
\end{equation}
where $r_{l} := m_1 + m_2 + \cdots + m_{k_l}$.
\end{proposition}

As we will see later in this paper (\eqref{eq:Man(8)} and \eqref{eq:Man(144)}, or Table~\ref{tab:m1m2}), 
$\manF(8) = \frac{13}{24}\log(8)$ and $\manF(144) = \frac{35}{128}\log(144)$. 
It can then be checked that \eqref{eq:Man(l)prod} is not valid for $l=8=F_3^3$ 
and $l = 144 = F_3^4 F_4^2$.
% Respectivamente, \eqref{eq:Man(l)prod} para $l=8=F_3^3$ y $l = 144 = F_3^4 F_4^2$ 
% dar\'ia $\frac{1}{24}\log(8)$ y $\frac{-5}{128}\log(144)$.

Proposition~\ref{prop:Man(l)prod} allows us to compute $\manF(l)$ for $l=F_{n}^m$ with $m \ge 1$ when $n \ne 3,6,12$ (i.e., $F_{n}\ne 2, 8, 144$), namely 
\[
\manF(F_{n}^m) = \frac{(-1)^{m-1}}{2^m} \log(F_{n}), \quad n \ne 3,6,12, \quad m \ge 1.
\]
The Fibonacci numbers $8$ and $144$ do not appear as basic factors in the decompositions~\eqref{eq:Fp}, but $F_3 = 2$ does. 
This proposition gives a closed expression for~$\manF(2^m)$:

\begin{proposition}
\label{prop:Man(2m)}
For $m \ge 1$, one has
\begin{align}
\label{eq:Man(2m)1}
  \manF(2^m) &= \left((-1)^{m+1} - 2^{-2m} (1+i\sqrt{7})^m 
  - 2^{-2m} (1-i\sqrt{7})^m\right) \log(2) \\
\label{eq:Man(2m)2}
  &= \left( (-1)^{m+1} - 2^{1-m/2} \cos(m \arctan\sqrt{7}) \right) \log(2) \\
\label{eq:Man(2m)3}
  &= \biggl( (-1)^{m+1} - 2^{1-2m} \sum_{k=0}^{\lfloor m/2 \rfloor} (-1)^k \binom{m}{2k} 7^k \biggr) \log(2).
\end{align}
\end{proposition}

For $l$ equal to $2^m$ times a Fibonacci number, we have the following:

\begin{proposition}
\label{prop:Man(2mFk)}
For $m \ge 1$, and a Fibonacci number $F_k \ge 3$, different from $8$ and $144$, one has
\begin{align*}
%\label{eq:Man(2mFk)1}
%\notag
  \manF(2^m F_k) &= \frac{1}{2^{2+m/2}} \left(2^{m/2}(-1)^m + \cos(m \arctan\sqrt{7}) 
  - \frac{1}{\sqrt{7}} \sin(m \arctan\sqrt{7}) \right) \log(2^mF_k) \\
%\label{eq:Man(2mFk)2}
%\notag
  &\sim \frac{(-1)^m}{4} \log(2^mF_k)
  \quad \text{when } m\to\infty.
\end{align*}
\end{proposition}

Notice that, as in \eqref{eq:Man(2m)3}, $\cos(m \arctan\sqrt{7})$ and $\frac{1}{\sqrt{7}} \sin(m \arctan\sqrt{7})$ can be rewritten as expressions without trigonometric functions or squared roots, indeed as rational numbers that depend on $m$, see~\eqref{eq:sqrt7}.

And, this is what happens when $l$ is $2^m$ times \emph{two} Fibonacci numbers:

\begin{proposition}
\label{prop:Man(2mFkFq)}
For $m \ge 1$, and two different Fibonacci numbers $F_k, F_q \ge 3$, also different from $8$ and $144$, one has
\begin{align*}
%\label{eq:Man(2mFkFq)1}
%\notag
  \manF(2^m F_kF_q) &= \frac{-1}{2^{5+m/2}} \biggl(2^{m/2}(-1)^m(2m+5) 
  + \frac{3}{7}(2m+7) \cos(m \arctan\sqrt{7}) \\
%\notag
  &\qquad\qquad\qquad\quad
  - \frac{1}{7\sqrt{7}} (14m+41) \sin(m \arctan\sqrt{7}) \biggr) \log(2^mF_kF_q) \\
%\label{eq:Man(2mFkFq)2}
%\notag
  &\sim \frac{(-1)^{m+1}}{16} \,m \log(2^mF_kF_q)
  \quad \text{when } m\to\infty.
\end{align*}
\end{proposition}

Propositions~\ref{prop:Man(2mFk)} and \ref{prop:Man(2mFkFq)} suggest that, for general values of $l$, it is not that easy to find closed expressions for $\manF(l)$. In Table~\ref{tab:m1m2} we can see the values of $\manF(2^{m_1} 3^{m_2})$ for $0 \le m_1 \le 8$ and $0 \le m_2 \le 5$, obtained with a computer algebra system (actually, we list the coefficients without the factor $\log(2^{m_1} 3^{m_2})$).

\begin{table}
%\centering
\[
%\newcommand{\vv}{\rule[-3pt]{0pt}{17pt}\relax\ignorespaces}
%\newcommand{\vv}{\rule[-7pt]{0pt}{19pt}\relax\ignorespaces}
\newcommand{\vv}{\rule[-1.5ex]{0pt}{4.25ex}\relax\ignorespaces}
\begin{array}{c|ccccccccc}
\vv & m_1=0 & m_1=1 & m_1=2 &
   m_1=3 & m_1=4 & m_1=5 &
   m_1=6 & m_1=7 & m_1=8 \\
\hline
\vv m_2=0 & - & \frac{1}{2} & -\frac{1}{8} &
   \frac{13}{24} & -\frac{17}{64} & \frac{21}{160} &
   -\frac{73}{384} & \frac{141}{896} & -\frac{225}{2048} \\
\vv m_2=1 & \frac{1}{2} & -\frac{1}{4} & \frac{1}{8}
   & -\frac{5}{16} & \frac{9}{32} & -\frac{13}{64} &
   \frac{33}{128} & -\frac{69}{256} & \frac{121}{512} \\
\vv m_2=2 & -\frac{1}{8} & \frac{1}{8} &
   -\frac{3}{32} & \frac{3}{16} & \frac{35}{128} &
   -\frac{5}{128} & -\frac{71}{512} & 0 & -\frac{81}{2048} \\
\vv m_2=3 & \frac{1}{24} & -\frac{1}{16} &
   \frac{1}{16} & -\frac{11}{96} & -\frac{11}{128} &
   \frac{17}{256} & \frac{19}{384} & \frac{9}{128} &
   -\frac{209}{2048} \\
\vv m_2=4 & -\frac{1}{64} & \frac{1}{32} &
   -\frac{5}{128} & \frac{9}{128} & \frac{13}{1024} &
   -\frac{11}{256} & -\frac{9}{1024} & -\frac{77}{1024} &
   \frac{475}{16384} \\
\vv m_2=5 & \frac{1}{160} & -\frac{1}{64} &
   \frac{3}{128} & -\frac{11}{256} & \frac{3}{256} &
   \frac{47}{2560} & -\frac{3}{1024} & \frac{123}{2048} &
   -\frac{205}{8192} \\
\hline
\end{array}
\]
\caption{The value of $\manF(2^{m_1}3^{m_2})/\log(2^{m_1}3^{m_2})$ for small values of $m_1$ and $m_2$.}
\label{tab:m1m2}
\end{table}

For general values of $l$, we can find some bounds. Note that in the following proposition we have introduced two numbers $p,q$ and in fact $p=1$; the notation is intended to allow speculation about what the effect of changing the values of $p,q$ would be.

\begin{proposition}
\label{prop:Mcg}
Let $\manF(l)$ be the arithmetical function defined in \eqref{eq:producto},
and
\[
l = 2^{m_1}3^{m_2}F_{n_3}^{m_3} \cdots F_{n_{k_l}}^{m_{k_l}} \in \Fp
\]
(factorization as in \eqref{eq:Fp}, which is unique).
Then, for a constant $C$ independent of $l$, one has
\begin{equation}
\label{eq:Mcg}
|\manF(l)| \leq 
\frac{C}{2^{r_l-pm_1-qm_2}} \frac{(r_l-1)!}{m_1!\,m_2!\,m_3! \cdots m_{k_l}!} \log(l),
\end{equation}
where $r_{l} := m_1+m_2+m_3+ \cdots +m_{k_l}$ (notice that $m_1$ and $m_2$ can be $=0$) with
\begin{equation}
\label{eq:pq022}
p=1 \quad\text{ and }\quad 
q = \log_2 \bigg( \frac{1}{2} + \frac{3}{10} \sqrt{5} \bigg)
= \log\bigg( \frac{1}{2} + \frac{3}{10} \sqrt{5} \bigg)/\log(2) = 0.22752\dots
\end{equation}
\end{proposition}

This is the final result of the paper:

\begin{theorem}
\label{theo:sigma0}
The Dirichet series \eqref{eq:ManFp} diverges for
\begin{equation}
\label{eq:diverge}
\Re(s) < \frac{7\log(7/2)}{\log(3\cdot 5\cdot 13\cdot 21\cdot 34\cdot 55\cdot 89)} = 0.431141\dots
\end{equation}
and converges absolutely for $\Re(s) > \sigma_0$, where $\sigma_0$
is the unique positive real number that satisfies $\varphi^{*}(s)=1$, with
\[
\varphi^{*}(s) :=
\frac{1}{2^{1-p}2^s} + \frac{1}{2^{1-q}3^s} + \frac{1}{2} \varphi(s) - 1
- \frac{1}{2 \cdot 2^s} - \frac{1}{2 \cdot 3^s} - \frac{1}{2 \cdot 8^s} - \frac{1}{2 \cdot 144^s}.
\]
With $p$ and $q$ as in \eqref{eq:pq022}, its approximate value is $\sigma_0 = 0.905556\dots$
\end{theorem}

%----------
\section{Expressions for the function $\manF(l)$: Proofs of Propositions~\ref{prop:Man(l)=0} up to \ref{prop:Man(2mFkFq)}}
%----------

%Observe that it is not necessary to include in the definition of $\calF%
%_{\p}$ the numbers $F_{6}=8$ and $F_{12}=144$ because already
%belong to $\Fp$: $8=2^{3}$, $144=2^{4}3^{2}$, so they
%are a power and a product of powers of the Fibonacci numbers $F_3=2$ and $%
%F_{4}=3$. Furthermore, it is important to stress that $F_{6}$ and $F_{12}$
%are the only Fibonacci numbers expressible as a power or as a product of
%powers of Fibonacci numbers \cite{BuLuMiSi, Car}.

%The convergence of $\sum_{n=1}^{\infty }F_{n}^{-s}$ on $\Re s>0$, guarantees
%the validity of the formula 
%\[
%\varphi'(s) = -\sum_{n=1}^{\infty }\left( \log F_{n}\right) F_{n}^{-s},
%\quad \Re s>0,
%\]
%and then, since $\varphi(s)\neq 0$ on $H_{\eta }$, \eqref{eq:ManFib} formally satisfies 
%\[
%\frac{\varphi'(s)}{\varphi(s)} 
%= -\sum_{l=1}^{\infty }\manF(l)l^{-s}, \quad s\in H_{\eta }. 
%\]
%Therefore 
%\[
%\sum_{n=1}^{\infty }\left( \log F_{n}\right) F_{n}^{-s}
%= \sum_{n=1}^{\infty} F_{n}^{-s} \sum_{l=1}^{\infty} \manF(l)l^{-s}, 
%\quad s\in H_{\eta }. 
%\]

From \eqref{eq:producto}, the coefficients $\manF(l)$ can be calculated by means of the
formal relation 
\[
\begin{aligned}
&\left( (\log 2) 2^{-s} + (\log 3) 3^{-s} + (\log 5) 5^{-s} + \cdots \right) 
= \left( 1+1+2^{-s}+3^{-s}+5^{-s}+\cdots\right) 
\\
&\qquad\qquad\qquad \times \left( \manF(1) + \manF(2)2^{-s} + \manF(3)3^{-s}
+ \manF(4)4^{-s} + \manF(5)5^{-s} + \manF(6)6^{-s}+\cdots\right) .
\end{aligned}
\]
It is clear that $\manF(1)=0$, so we may rewrite this as
\begin{equation}
\label{eq:Mansimple}
\begin{aligned}
&\left( (\log 2) 2^{-s} + (\log 3) 3^{-s} + (\log 5) 5^{-s} + \cdots\right) 
= \left( 2+2^{-s}+3^{-s}+5^{-s}+\cdots\right) 
\\
&\qquad\qquad\qquad \times \left( \manF(2)2^{-s} + \manF(3)3^{-s} + \manF(4)4^{-s}
+ \manF(5)5^{-s} + \manF(6)6^{-s} + \cdots \right) . 
\end{aligned}
\end{equation}

Consider an integer $l\ge 2$ and let $F_{d_1},\dots, F_{d_{p_l}}$ be the proper Fibonacci divisors of $l$ (as usual, a proper divisor of a number is a divisor that 
is not equal to the number itself). Then, directly from \eqref{eq:Mansimple}, 
we deduce that $\manF(l)$ can be recursively determined by means of the formula
\begin{equation}
\label{eq:Man(l)rec}
\manF(l) = 
\begin{cases}
-\frac{1}{2} (\manF(j_1) + \cdots + \manF(j_{p_l}) - \log l),
& \text{ if } l\in \calF, \quad p_{l} \geq 1, \\ 
-\frac{1}{2}(\manF(j_1) + \cdots + \manF(j_{p_l})),
& \text{ if } l\notin \calF, \quad p_{l} \geq 1, \\ 
\frac{1}{2} (\log l),
& \text{ if } l\in \calF, \quad p_{l} = 0, \\ 
0,
& \text{ if } l\notin \calF, \quad p_{l} = 0,
\end{cases}
\end{equation}
where
\[
j_1 = \frac{l}{F_{d_1}},\ \dots,\
j_{p_l} = \frac{l}{F_{d_{p_l}}}, \qquad p_{l} \geq 1.
\]

Let us illustrate the use of \eqref{eq:Man(l)rec} with a couple of examples that have special interest. For instance, $8\in \calF$ has $2\in \calF$ as its unique Fibonacci
proper divisor. Therefore, $p_{8}=1$ and, by \eqref{eq:Man(l)rec}, one has 
$\manF(8) = -\frac{1}{2}(\manF(4)-\log(8))$. Now, since $4\notin \calF$
and $2$ is a Fibonacci proper divisor of $4$, $p_{4}=1$ and then, by \eqref{eq:Man(l)rec}, 
$\manF(4) = -\frac{1}{2} \manF(2)$. Thus we get 
$\manF(8) = \frac{1}{2} \left(\log(8) + \frac{1}{2} \manF(2)\right)$. 
Finally, $2\in \calF$ but $2$ has no
Fibonacci proper divisor greater than $1$, so $p_2=0$ and then 
$\manF(2) = \frac{1}{2} \log(2)$. Therefore 
\begin{equation}
\label{eq:Man(8)}
\manF(8) = \frac{13}{24} \log(8). 
\end{equation}
A more cumbersome example is $144 = 2^4\cdot3^2 \in \calF$. Its Fibonacci proper divisors are $2$, $3$ and $8$. By \eqref{eq:Man(l)rec}, we have
\[
\manF(144) = -\frac{1}{2}(\manF(72)+\manF(48)+\manF(18)-\log(144)),
\]
and again we use \eqref{eq:Man(l)rec} to compute $\manF(72)$, $\manF(48)$, $\manF(18)$,
and so on. This recursive process finally yields
\begin{equation}
\label{eq:Man(144)}
\manF(144) = \frac{35}{128} \log(144). 
\end{equation}
This method can be used to find all the values that appear in Table~\ref{tab:m1m2}.
Of course, it can be implemented in any computer algebra system.

%% Ejemplo que se podr\'ia usar: $l=(21^2)(34^4)$
\begin{remark}
It is important to observe the following regarding the practical application of \eqref{eq:Man(l)rec}. Let $l \in \Fp$, with factorization
\[
l = F_{n_1}^{m_1}F_{n_2}^{m_2} \cdots F_{n_{k_l}}^{m_{k_l}}
\]
as in \eqref{eq:Fp} (i.e., with $F_{n_i} \ne 8,144$), which is unique.
If $F_d > 1$ is a Fibonacci number that divides $l$, we cannot ensure that $F_d$ is one of the $F_{n_i}$ that appear in the above decomposition. This happens often, due to the well-known property $\gcd(F_a,F_b) = F_{\gcd(a,b)}$ of the Fibonacci numbers. But, if a proper divisor $F_d$ is not $8$, $144$, or one of the $F_{n_i}$, then in fact $l/F_d \notin \Fp$:  indeed, if $l' = l/F_d \in \Fp$, then let
\[
  \frac{l}{F_d} = l' = F_{n_1'}^{m_1'} F_{n_2'}^{m_2'} \cdots F_{n_{k_{l'}}'}^{m_{k_{l'}}'}
\]
be its factorization. This leads to two different factorizations of $l$,
\[
 l = F_{n_1}^{m_1}F_{n_2}^{m_2} \cdots F_{n_{k_l}}^{m_{k_l}}
 = F_d \cdot F_{n_1'}^{m_1'} F_{n_2'}^{m_2'} \cdots F_{n_{k_{l'}}'}^{m_{k_{l'}}'},
\]
which contradicts uniqueness.

Moreover, also by unique factorization, $l/8 \notin \Fp$ if $F_d = 8$ is a proper divisor of $l$ and the factorization of $l$ does not contain $F_{n_1}=2$ with $m_1\ge 3$; and $l/144 \notin \Fp$ if $F_d = 144$ is a proper divisor of $l$ and $l$ does not have the factor $F_{n_1}=2$ with $m_1\ge 4$ and $F_{n_2}=3$ with $m_2\ge 2$.
On the other hand, Proposition~\ref{prop:Man(l)=0} shows that $\manF(l) = 0$ when $l \notin \Fp$. Thus, after we prove this result, we do not need, in \eqref{eq:Man(l)rec}, to take all the $F_d$ that divide $l$, only the proper divisors that are $8$, $144$ (if $F_{n_1}=2$ with $m_1\ge 3$, or if $F_{n_1}=2$ with $m_1\ge 4$ and $F_{n_2}=3$ with $m_2\ge 2$, respectively), or one of the Fibonacci numbers $F_{n_i}$ that appear in the decomposition of $l$ in $\Fp$ because, otherwise, the corresponding summand $\manF(l/F_d)$ is null. This will be tacitly used many times in the forthcoming proofs. 
\end{remark}

Using \eqref{eq:Man(l)rec}, let us now prove Propositions~\ref{prop:Man(l)=0}, 
\ref{prop:Man(l)prod}, \ref{prop:Man(2m)}, \ref{prop:Man(2mFk)}, 
and~\ref{prop:Man(2mFkFq)}.

%----------
\subsection{Proof of Proposition~\ref{prop:Man(l)=0}}
%----------

%% PARECE QUE ESTA DEMO QUE SIGUE PROPORCIONA UNA DEMO PARA UN CASO MÁS GENERAL:
% Si empezamos con un subconjunto cualquiera F de enteros (en este caso los Fibonacci), 
% que contenga el 1, entonces la inversa de Dirichlet de la funci\'on caracter\'istica 1_F de F 
% est\'a soportada en el semigrupo generado por F, es decir, 
% se anula sobre números que no son productos de elementos de F.

%\begin{proof}[Proof of Proposition~\ref{prop:Man(l)=0}]
We will prove that $\manF(l) = 0$ when $l \notin \Fp$ by complete induction on $l$.
The first $l$ that does not belong to $\Fp$ is $l=7$, and $\manF(7)=0$ by \eqref{eq:Man(l)rec} (the case $l \notin \calF$ and $p_l=0$). Let us now assume that $\manF(l') = 0$ for any $l'<l$ such that  $l' \notin \Fp$. We want to prove that also $\manF(l) = 0$. Let us write \eqref{eq:Mansimple} as
\begin{equation}
\label{eq:ManpD}
\sum_{k=2}^\infty \log(F_k) F_k^{-s} 
= \left(\sum_{n=1}^\infty a_{n} n^{-s}\right) \left(\sum_{n=1}^\infty \manF(n) n^{-s}\right)
= \sum_{n=1}^\infty \left( \sum_{k \mid n} a_k \manF(n/k) \right) n^{-s},
\end{equation}
where $a_1=2$, $a_k=1$ if $k=F_j$ for some $j\ge3$ (starting with $k=F_3=2$) and $a_k=0$ otherwise.
Comparing the coefficients of $l^{-s}$ in~\eqref{eq:ManpD}, the coefficient on the left-hand side is $0$ because $l \notin \Fp$, so $l$ is not a Fibonacci number; then if on the right-hand side we separate the summand corresponding to $k=1$, we have
\[
0 = a_1 \manF(l) + \sum_{k \mid l,\, k > 1} a_k \manF(l/k).
\]
Here, $a_k=0$ except when $k=F_j$. And, since $l \notin \Fp$, also $l/F_j \notin \Fp$, so, by the inductive hypothesis, $\manF(l/k) = 0$ for every $k=F_j>1$. Consequently, $\manF(l)=0$, and the proof is concluded.
%\end{proof}

%% En lugar de aqu\'i, esto tambi\'en se pod\'ia poner, sin llamarlo ``Remark'', 
%% en la introducci\'on, justo después de la Proposition~\ref{prop:Man(l)=0}
\begin{remark}
This raises the question of whether we can have $\manF(l) = 0$ with $l \in \Fp$. 
In light of \eqref{eq:Man(l)prod}, this can only happen when $l \in \Fp$ does not satisfy the hypotheses of Proposition~\ref{prop:Man(l)prod}. In fact $\manF(2^{7}3^{2}) = 0$ is the only $l$ in the range $1<l\le10^8$ such that $l \in \Fp$ and $\manF(l) = 0$.
\end{remark}
%% No parece f\'acil llegar m\'as lejos, salvo a trozos; tarda ya mucho
%% y Mathematica se queda a veces sin memoria, y hay que reiniciarlo y seguir

%----------
\subsection{Proof of Proposition~\ref{prop:Man(l)prod}}
%----------

%\begin{proof}[Proof of Proposition~\ref{prop:Man(l)prod}]
If $l$ is a number of $\Fp$ with $r_{l}=1$ that fulfills the hypothesis, 
$l$ is necessarily a Fibonacci number distinct from $F_{6}=8$ and $F_{12}=144$
(observe that $r_{8}=3$, $r_{144}=6$). Then, directly by \eqref{eq:Man(l)rec} one has 
\[
\manF(l) = \frac{1}{2} \log(l), 
\]
so formula \eqref{eq:Man(l)prod} is true if $r_{l}=1$. 

To prove the validity of \eqref{eq:Man(l)prod} in general we use complete induction on $r_{l}$. Let $l = F_{n_1}^{m_1} F_{n_2}^{m_2} \cdots F_{n_{k_l}}^{m_{k_l}}$ be an element of $\Fp$ (with $F_{n_1} \neq 2$ or $F_{n_1}=2$ and $m_1<3$) having $r_{l}>1$ and assume that \eqref{eq:Man(l)prod} is true for any 
\[
l' = F_{n_1'}^{m_1'} F_{n_2'}^{m_2'} 
\cdots F_{n_{k_{l'}}'}^{m_{k_{l'}}'} \in \Fp
\]
with 
$r_{l'} := m_1' + m_2' + \cdots + m_{k_{l'}}' < r_{l}$
(and such that $F_{n_1'} \neq 2$ or $F_{n_1'}=2$ and $m_1'<3$).
By \cite{BuLuMiSi, Car}, $l\notin \calF$ and $l/F_j \notin \Fp$ 
if $F_j$ is a Fibonacci proper divisor of $l$ that is not one of the $F_{n_{i}}$.
Then, by applying \eqref{eq:Man(l)rec} and Proposition~\ref{prop:Man(l)=0}, we have 
\begin{equation}
\label{eq:Man(l)demo}
\manF(l) = -\frac{1}{2} \left(  \manF(F_{n_1}^{m_1-1} \cdots F_{n_{k_l}}^{m_{k_l}})
+ \manF(F_{n_1}^{m_1}F_{n_2}^{m_2-1} \cdots F_{n_{k_l}}^{m_{k_l}})
+ \cdots + \manF(F_{n_1}^{m_1} \cdots F_{n_{k_l}}^{m_{k_l}-1}) \right) .
\end{equation}
Since the sum of exponents of each summand of the right-hand side of \eqref{eq:Man(l)demo}
is $r_l-1$, the induction hypothesis applies and we obtain
\begin{align*}
\manF(l) &= -\frac{1}{2}(-1)^{r_l-2} \,\frac{1}{2^{r_l-1}}(r_l-2)! \\
& \qquad\quad \times
\left( \frac{1}{(m_1-1)! \cdots m_{k_l}!} \log (l/F_{n_1}) + \cdots 
+ \frac{1}{m_1! \cdots (m_{k_l}-1)!} \log (l/F_{n_{k_l}}) \right) \\
&= (-1)^{r_l-1}\frac{1}{2^{r_l}} \frac{(r_l-1)!}{m_1! \cdots m_{k_l}!} \log(l), 
\end{align*}
which proves~\eqref{eq:Man(l)prod}. 
%\end{proof}

%----------
\subsection{Proof of Proposition~\ref{prop:Man(2m)}}
%----------

%\begin{proof}[Proof of Proposition~\ref{prop:Man(2m)}]
We have already seen that $\manF(2) = \frac{1}{2} \log(2)$, 
$\manF(2^2) = -\frac{1}{4} \log(2)$ and $\manF(2^3) = \frac{13}{8} \log(2)$.
For $m \ge 4$, the Fibonacci proper divisors of $l= 2^m$ are $2$ and $8$, 
so \eqref{eq:Man(l)rec} gives
\[
\manF(2^m) = -\frac{1}{2} \left( \manF(2^{m-1}) + \manF(2^{m-3}) \right).
\]
For simplicity, let us write $a_m = \manF(2^m)$, which satisfies the linear homogeneous recurrence 
\begin{equation}
\label{eq:am}
a_m + \frac{1}{2} a_{m-1} + \frac{1}{2} a_{m-3} = 0,
\end{equation}
whose characteristic equation is $\lambda^3 + \frac{1}{2} \lambda^2 + \frac{1}{2} = 0$.
The roots of the characteristic equation are 
\[
\lambda_1 = -1 \quad\text{and}\quad
\lambda_{2,3} = \frac{1}{4}\left(1 \pm i\sqrt{7}\right) 
= \frac{1}{\sqrt{2}} \left(\cos(\arctan\sqrt{7}) \pm i \sin(\arctan\sqrt{7}) \right).
\] 
Thus, the solution of \eqref{eq:am} can be written as
\[
a_m = A \lambda_1^m + B \lambda_2^m + C \lambda_3^m
\]
or also as
\[
a_m = A \lambda_1^m + B' \left(\frac{1}{\sqrt{2}}\right)^m \cos(m\arctan\sqrt{7}) 
+ C' \left(\frac{1}{\sqrt{2}}\right)^m \sin(m\arctan\sqrt{7}),
\]
where the constants $A,B,C$ or $A,B',C'$ must be chosen to satisfy the initial values
$a_1 = \frac{1}{2} \log(2)$, $a_2 = -\frac{1}{4} \log(2)$ and $a_3 = \frac{13}{8} \log(2)$.
It is a simple task to find the values of these constants, and thus we get \eqref{eq:Man(2m)1} and~\eqref{eq:Man(2m)2}. 
Finally, 
% the sum of the binomial expansions  
% $(1 \pm i \sqrt{7})^m = \sum_{k=0}^{m} \binom{m}{k} (\pm i)^k (\sqrt{7})^k$ is
\[
(1 + i \sqrt{7})^m + (1 - i \sqrt{7})^m
= \sum_{k=0}^{m} \binom{m}{k} i^k (\sqrt{7})^k + \sum_{k=0}^{m} \binom{m}{k} (-i)^k (\sqrt{7})^k
= 2 \sum_{k=0}^{\lfloor m/2 \rfloor} \binom{m}{2k} (-1)^k 7^k,
\]
which becomes \eqref{eq:Man(2m)3}.
%\end{proof}

%----------
\subsection{Proof of Proposition~\ref{prop:Man(2mFk)}}
%----------

%\begin{proof}[Proof of Proposition~\ref{prop:Man(2mFk)}]
By using the recurrence \eqref{eq:Man(l)rec}, we easily get
\[
\manF(2F_k) = -\frac{1}{4} \log(2F_k),
\quad
\manF(2^2F_k) = \frac{1}{8} \log(2^2F_k),
\quad
\manF(2^3F_k) = -\frac{5}{16} \log(2^3F_k).
\]
Also, for $m \ge 3$,
\[
\manF(2^mF_k) = -\frac{1}{2} \left( \manF(2^{m-1}F_k)
+ \manF(2^{m-3}F_k) + \manF(2^{m}) \right).
\]
For simplicity, let $a_m = \manF(2^mF_k)$, and $b(m) = -\frac{1}{2}\manF(2^{m})$, which is given by~\eqref{eq:Man(2m)1}. Then we have the linear recurrence
\begin{equation*}
%\label{eq:amFm}
\begin{gathered}
a_m + \frac{1}{2} a_{m-1} + \frac{1}{2} a_{m-3} = b(m),
\\
a_1 = -\frac{1}{4} \log(2F_k),
\quad
a_2 = \frac{1}{8} \log(2^2F_k),
\quad
a_3 = -\frac{5}{16} \log(2^3F_k),
\end{gathered}
\end{equation*}
with $b(m) = c_1 p_1^m + c_2 p_2^m + c_3 p_3^m$
for certain constants $c_1,c_2,c_3$ and $p_1,p_2,p_3$.
Solving this recurrence by standard methods, we get
\[
a_m = \frac{1}{2^{2+m/2}} \left(2^{m/2}(-1)^m + \cos(m \arctan\sqrt{7}) 
  - \frac{1}{\sqrt{7}} \sin(m \arctan\sqrt{7}) \right) \log(2^mF_k).
\]
Finally, it is clear from this expression that 
$a_m \sim \frac{(-1)^m}{4} \log(2^mF_k)$ when $m \to \infty$.
%\end{proof}

\begin{remark}
Let us see how the expressions $\cos(m \arctan\sqrt{7})$ and $\frac{1}{\sqrt{7}} \sin(m \arctan\sqrt{7})$ that appear, among other places, in Proposition~\ref{prop:Man(2mFk)}, can be rewritten without trigonometric functions or square roots. We start with De Moivre's formula
\[
\cos (m\theta) + i \sin(m\theta) = (\cos(\theta) + i \sin(\theta))^{m}.
\]
Applying the binomial formula and equating real and imaginary parts, we obtain
%\[
%  \cos (m \theta) = \sum_{r=0}^{2r \leq m} (-1)^{r} 
%  \binom{m}{2r} \cos^{m-2r}(\theta) \sin^{2r}(\theta)
%  = \cos^{m}(\theta) \sum_{r=0}^{2r \leq m} (-1)^{r} 
%  \binom{m}{2r} \tan^{2r}(\theta)
%\]
\[
  \cos (m \theta) = \sum_{r=0}^{\lfloor m/2\rfloor} (-1)^{r} 
  \binom{m}{2r} \cos^{m-2r}(\theta) \sin^{2r}(\theta)
  = \cos^{m}(\theta) \sum_{r=0}^{\lfloor m/2\rfloor} (-1)^{r} 
  \binom{m}{2r} \tan^{2r}(\theta)
\]
and
%\[
%  \sin(m \theta) = \sum_{r=0}^{2r+1 \leq m} (-1)^{r} 
%  \binom{m}{2r+1} \cos^{m-2r-1}(\theta) \sin^{2r+1}(\theta)
%  = \cos^{m}(\theta) \sum_{r=0}^{2r+1 \leq m} (-1)^{r} 
%  \binom{m}{2r+1} \tan^{2r+1}(\theta).
%\]
\begin{align*}
  \sin(m \theta) &= \sum_{r=0}^{\lfloor(m-1)/2\rfloor} (-1)^{r} 
  \binom{m}{2r+1} \cos^{m-2r-1}(\theta) \sin^{2r+1}(\theta)
  \\&= \cos^{m}(\theta) \sum_{r=0}^{\lfloor(m-1)/2\rfloor} (-1)^{r} 
  \binom{m}{2r+1} \tan^{2r+1}(\theta).
\end{align*}
Now, using $\cos^2(\theta) = 1/(1+\tan^2(\theta))$, we get
\begin{align*}
  \cos (m \theta) &= \frac{1}{(1+\tan^{2}(\theta))^{m/2}} 
  \sum_{r=0}^{\lfloor m/2\rfloor} (-1)^{r} 
  \binom{m}{2r} \tan^{2r}(\theta),
\\
  \sin(m \theta) &= \frac{1}{(1+\tan^{2}(\theta))^{m/2}} 
  \sum_{r=0}^{\lfloor(m-1)/2\rfloor} (-1)^{r} 
  \binom{m}{2r+1} \tan^{2r+1}(\theta).
\end{align*}
Finally, taking $\theta = \arctan(\sqrt{7})$, we obtain
\begin{equation}
\label{eq:sqrt7}
\begin{aligned}
  \cos (m \arctan(\sqrt{7})) &= \frac{1}{8^{m/2}}
  \sum_{r=0}^{\lfloor m/2\rfloor} (-1)^{r} 7^r \binom{m}{2r},
\\
  \frac{1}{\sqrt{7}} \sin(m \arctan(\sqrt{7})) &= \frac{1}{8^{m/2}}
  \sum_{r=0}^{\lfloor(m-1)/2\rfloor} (-1)^{r} 7^r \binom{m}{2r+1}.
\end{aligned}
\end{equation}
\end{remark}

%----------
\subsection{Proof of Proposition~\ref{prop:Man(2mFkFq)}}
%----------

%\begin{proof}[Proof of Proposition~\ref{prop:Man(2mFkFq)}]
By using the recurrence \eqref{eq:Man(l)rec}, we easily get
\[
\manF(2F_kF_q) = \frac{1}{4} \log(2F_kF_q),
\quad
\manF(2^2F_kF_q) = -\frac{3}{16} \log(2^2F_kF_q),
\quad
\manF(2^3F_kF_q) = \frac{3}{8} \log(2^3F_kF_q).
\]
In addition, for $m \ge 3$,
\[
\manF(2^mF_kF_q) = -\frac{1}{2} \left( \manF(2^{m-1}F_kF_q)
+ \manF(2^{m-3}F_kF_q) + \manF(2^{m}F_k) + \manF(2^{m}F_q) \right).
\]
Let $a_m = \manF(2^mF_kF_q)$ and $b_k(m) = -\frac{1}{2}\manF(2^{m}F_k)$, given in Proposition~\ref{prop:Man(2mFk)}. Then we have the linear recurrence
\begin{equation*}
\begin{gathered}
a_m + \frac{1}{2} a_{m-1} + \frac{1}{2} a_{m-3} = b_k(m) + b_q(m),
\\
a_1 = \frac{1}{4} \log(2F_kF_q),
\quad
a_2 = -\frac{3}{16} \log(2^2F_kF_q),
\quad
a_3 = \frac{5}{8} \log(2^3F_kF_q).
\end{gathered}
\end{equation*}
Solving the recurrence, we get the value for $a_m$. In this case solving the recurrence is a cumbersome process (it can be done with the help of computer algebra software). But, after the expression for $a_m$ has been found, checking that it satisfies the recurrence is not difficult. Finally, the asymptotic formula for $a_m$ follows immediately from its expression.
%\end{proof}

%----------
\subsection{Some additional expressions}
%----------

% Decir que esto es \'acil de observar en la pr\'actica aunque no incluimos una demostración formal.

It is easy to verify on a case by case basis that $\manF(l)$ is always a rational number multiplied by $\log(l)$. In addition, we observe that the value of $\manF(l)$, with $l \in \Fp$ factored as in \eqref{eq:Fp}, exhibits symmetry with respect to those factors different from $2$ and $3$. Namely, for $m_1 \ge 0$, $m_2 \ge 0$, and different Fibonacci numbers $F_{n_j}\ge 5$, also different from $8$ and $144$, let us denote
\begin{equation*}
%\label{eq:Man(2m3nprods)}
\widehat{\manF}(m_1,m_2,J) := 
\frac{\manF\Big(2^{m_1} 3^{m_2} \prod_{j=1}^{J} F_{n_j}\Big)}
  {\log\Big(2^{m_1} 3^{m_2} \prod_{j=1}^{J} F_{n_j}\Big)},
\qquad\text{with } m_1+m_2+J>0.
\end{equation*}
%% De todas formas, el caso J=0 es trivial
Then, for $l=2^{m_1} 3^{m_2} F_{n_3}^{m_3} F_{n_4}^{m_4} \cdots F_{n_{k_l}}^{m_{k_l}} > 1$ (factored uniquely as in \eqref{eq:Fp}), one has
\[
  \frac{\manF(2^{m_1} 3^{m_2} F_{n_3}^{m_3} \cdots F_{n_{k_l}}^{m_{k_l}})}
    {\log(2^{m_1} 3^{m_2} F_{n_3}^{m_3} \cdots F_{n_{k_l}}^{m_{k_l}})}
  = \frac{1}{m_3!\,m_4!\cdots m_{k_l}!} 
  \, \widehat{\manF}(m_1,m_2,m_3+m_4+\cdots+m_{k_l}).
\]

%----------
\section{Bounds for the function $\manF(l)$: Proof of Proposition~\ref{prop:Mcg}}
%----------

Clearly, the bound \eqref{eq:Mcg} in Proposition~\ref{prop:Mcg} holds for $l$ satisfying the hypotheses of Proposition~\ref{prop:Man(l)prod} (in that case, $p$ and $q$ do not play any role). For $l \in \Fp$ in general, the idea is to give an inductive proof using the recurrence~\eqref{eq:Man(l)rec}. If $l = 2^{m_1}3^{m_2}F_{n_3}^{m_3} \cdots F_{n_{k_l}}^{m_{k_l}}$, then $p$ and $q$ in \eqref{eq:Mcg} can be considered as functions
\[
  p := p(m_1,r_l),\quad q := q(m_2,r_l),
\]
with values in the interval $[0,1]$, which must be chosen in such a way that~\eqref{eq:Mcg} holds.

As we shall see in the following proof, this is easy if we are satisfied with $p = q = 1$; however, this choice yields a poor bound. The smaller $p$ and $q$ are, the better the bound is for use in the proof of Theorem~\ref{theo:sigma0}, but the effort that must be made to obtain these improved bounds is proportionally greater. In addition, note that if we want to use \eqref{eq:Mcg} to prove Theorem~\ref{theo:sigma0}, then we need global bounds where $p$ and $q$ are independent from $m_1$, $m_2$ and~$r_l$.

In light of Propositions \ref{prop:Man(2m)}, \ref{prop:Man(2mFk)} and~\ref{prop:Man(2mFkFq)}, the bound \eqref{eq:Mcg} cannot hold for $p=q=0$; furthermore, for such a bound to hold for all $l \in \Fp$, it must necessarily be the case that $p = 1$. This does not exclude the possibility that bounds of a different type could exist. For example, numerical experiments seem to point to the existence of a bound of the form
\[
|\manF(l)| \leq 
\frac{C}{2^{r_l-m_1+t\min(m_1,m_2)}}
\dfrac{(r_l-1)!}{m_1!\,m_2!\,m_3! \cdots m_{k_l}!} \log(l),
\]
where $t$ is very close to~$1$. However, we have not been able to prove such a bound, and even if we had, it does not seem straightforward to use it to improve Theorem~\ref{theo:sigma0}).

In fact, when we bound $|\manF(l)|$ via \eqref{eq:Man(l)rec}, we are taking absolute values in all the summands, which of course discards possible cancelation due to alternating signs in $\manF(l)$. Proposition~\ref{prop:Man(l)prod} suggests that $\manF(l)$ changes sign each time that a Fibonacci proper divisor is extracted from $l$; however, this is not always the case, as one can see in Table~\ref{tab:m1m2}, which makes any proof based on cancelation difficult, since one needs to take advantage of a phenomenon which happens often, but not always, and not in a controlled or localized manner.

Hence, we are convinced that $|\manF(l)|$ has better bounds than those which we will discuss here, as supported by numerical experiments, although obtaining proofs does not seem easy.

Let us then turn to the proof of Proposition~\ref{prop:Mcg}. We have set it up so as to exhibit the difficulties that arise if we want $p, q \in [0,1]$ to be small. In Remark~\ref{rem:m1} we discuss the possibility of finding bounds for $|\manF(l)|$ with $p<1$, by dropping the requirement that they hold for all $l$ and instead impose restrictions, such as $r_{l} \geq 2m_{1}$.

\begin{proof}[Proof of Proposition~\ref{prop:Mcg}]
We want to prove that, for a certain constant $C_{p,q}$ (that depends only on $p$ and $q$), 
\begin{equation}
\label{eq:McgC}
|\manF(l)| \leq 
\frac{C_{p,q}}{2^{r_l-pm_1-qm_2}} \frac{(r_l-1)!}{m_1!\,m_2!\,m_3! \cdots m_{k_l}!} \log(l),
\end{equation}
where $r_{l} := m_1+m_2+m_3+ \cdots +m_{k_l}$ (note that $m_1$ and $m_2$ can be~$=0$).

To deal with the case $l=2^{m_1}$ (recall Proposition~\ref{prop:Man(2m)}), it is enough to take $p=1$ and $C_{p,q}$ a little larger than $1$, for example $C_{p,q}=7/4$ (note that if the bound holds for $q=0$, then it also holds for any $q\in[0,1]$);
for the moment, though, let us forget about $l=2^{m_1}$. We will see that even when we look for a bound not covering that case, we also need to take $p = 1$.

Thus, we are going to prove \eqref{eq:McgC} for 
\begin{equation*}
l = 2^{m_1}3^{m_2}F_{n_3}^{m_3} \cdots F_{n_{k_l}}^{m_{k_l}} \in \Fp
\end{equation*}
(factorization as usual as in \eqref{eq:Fp}) by induction on $m_1$. We will not bother with checking in detail what the exact value of $C_{p,q}$ is; when the argument requires discarding a finite number of values, it might be necessary to make $C_{p,q}$ bigger. In fact from now on we work with a global value of $C = 7/4$.

Proposition~\ref{prop:Man(l)prod} shows that \eqref{eq:McgC} is true when $m_1 \in \{0,1,2\}$ (in this case, the bound is valid for $p=q=0$, hence also for any $p,q\in[0,1]$), so it is enough to see what happens when $m_1 \ge 3$.
Indeed, this is the starting point of the inductive process.

In order to avoid dealing with many separate cases, and to allow us to express summands which do not actually form part of some expressions, we will use the following notation:
\[
  \IF_{\mathrm{condition}}(a) =
  \begin{cases}
  a & \text{if ``condition'' is true},\\
  0 & \text{if ``condition'' is false}.
  \end{cases}
\]
Note that, in the following arguments, we always bound $\IF_{\mathrm{condition}}(a)$ by a positive constant, so the bound can also be used when $\IF_{\mathrm{condition}}(a) = 0$. We will use two different ``conditions'', namely
\[
(m_2 \ge 1)  \quad\text{and}\quad (m_1 \ge 4 \text{ and } m_2 \ge 2)
\]
which, for simplicity, will be denoted by $\IF_3$ and $\IF_{144}$.
Observe that if they do not hold, the corresponding summand $\manF(l/3)$ or $\manF(l/144)$ in \eqref{eq:Man(l)rec} is absent. 

%% Para estudiar lo que sigue, los $F_i$ que no sean $2$ y $3$ pr\'acticamente no pintan nada. 
%% Es como si no estuviesen.

Under the inductive hypothesis, we assume that the bound \eqref{eq:McgC} holds for all $l'$ whose corresponding $m_1'$ is less than $m_1$. By the recursive formula \eqref{eq:Man(l)rec}, we have 
\begin{align*}
%\notag
%\label{eq:MFl-1}
|\manF(l)| &\leq \frac{1}{2} \left| \manF(2^{m_1-1}3^{m_2}F_{n_3}^{m_3} \cdots F_{n_{k_l}}^{m_{k_l}}) \right|
+ \frac{1}{2} \left| \manF(2^{m_1-3}3^{m_2}F_{n_3}^{m_3} \cdots F_{n_{k_l}}^{m_{k_l}})\right| 
\\
%\notag
&\quad+ \frac{1}{2} \IF_3 \left( \left| \manF(2^{m_1}3^{m_2-1}F_{n_3}^{m_3} 
  \cdots F_{n_{k_l}}^{m_{k_l}})\right| \right) 
+ \frac{1}{2} \IF_{144} \left( \left| \manF(2^{m_1-4}3^{m_2-2}F_{n_3}^{m_3} 
  \cdots F_{n_{k_l}}^{m_{k_l}})\right| \right) 
\\
%\label{eq:MFl-2}
%\notag
&\quad + \frac{1}{2} \left| \manF(2^{m_1}3^{m_2}F_{n_3}^{m_3-1} \cdots F_{n_{k_l}}^{m_{k_l}}) \right|
+ \cdots + \frac{1}{2} \left| \manF(2^{m_1}3^{m_2}F_{n_3}^{m_3} \cdots F_{n_{k_l}}^{m_{k_l}-1}) \right|
\\
&=: \frac{1}{2} S_1 + \frac{1}{2} S_2 + \frac{1}{2} S_3  
+ \frac{1}{2} S_4 + \frac{1}{2} S_5 + \frac{1}{2} S_6.
%\label{eq:men6}
\end{align*}
By induction,
\begin{align}
\notag
% &\left| \manF(2^{m_1-1}3^{m_2}F_{n_3}^{m_3} \cdots F_{n_{k_l}}^{m_{k_l}}) \right|
S_1 
&\le \frac{C}{2^{r_l-1-p(m_1-1)-qm_2}} \frac{(r_l-2)!}{(m_1-1)!\,m_2!\,m_3! \cdots m_{k_l}!} \log(l/2) \\
&= \frac{C}{2^{r_l-pm_1-qm_2}} \frac{(r_l-2)!}{m_1!\,m_2!\,m_3! \cdots m_{k_l}!} 
\, 2^{1-p} m_1 \log(l/2)
\label{eq:men6-1}
\end{align}
and again by induction,
\begin{align}
\notag
%&\left| \manF(2^{m_1-3}3^{m_2}F_{n_3}^{m_3} \cdots F_{n_{k_l}}^{m_{k_l}}) \right|
S_2 
&\le \frac{C}{2^{r_l-3-p(m_1-3)-qm_2}} \frac{(r_l-4)!}{(m_1-3)!\,m_2!\,m_3! \cdots m_{k_l}!} \log(l/8) \\
&= \frac{C}{2^{r_l-pm_1-qm_2}} \frac{(r_l-2)!}{m_1!\,m_2!\,m_3! \cdots m_{k_l}!}
\frac{m_1(m_1-1)(m_1-2)}{(r_l-2)(r_l-3)} \, 2^{3-3p} \log(l/8)
\label{eq:men6-2}
\end{align}
(if needed, we can use that $\log(l/8) \le \log(l/2)$).
By induction (noting the summand is absent if $\IF_3$ is false)
\begin{align}
\notag
%\left| \manF(2^{m_1}3^{m_2}F_{n_3}^{m_3} \cdots F_{n_{k_l}}^{m_{k_l}}) \right| 
S_3 
&\le \frac{C}{2^{r_l-1-pm_1-q(m_2-1)}} \frac{(r_l-2)!}{m_1!\,(m_2-1)!\,m_3! \cdots m_{k_l}!} \log(l/3) \\
&= \frac{C}{2^{r_l-pm_1-qm_2}} \frac{(r_l-2)!}{m_1!\,m_2!\,m_3! \cdots m_{k_l}!} 
\, 2^{1-q} m_2 \log(l/3).
\label{eq:men6-3}
\end{align}
By induction (noting that this summand is absent if $\IF_{144}$ is false), 
\begin{align}
\notag
%\left| \manF(2^{m_1-4}F_{n_3}^{m_3} \cdots F_{n_{k_l}}^{m_{k_l}})\right|
S_4 
&\le \frac{C}{2^{r_l-6-p(m_1-4)-q(m_2-2)}} \frac{(r_l-7)!}{(m_1-4)!\,(m_2-2)!\,m_3! \cdots m_{k_l}!} \log(l/144) \\
&= \frac{C}{2^{r_l-pm_1-qm_2}} \frac{(r_l-2)!}{m_1!\,m_2!\,m_3! \cdots m_{k_l}!} 
\frac{m_1(m_1-1)(m_1-2)(m_1-3)m_2(m_2-1)}{(r_l-2)(r_l-3)(r_l-4)(r_l-5)(r_l-6)} \, 2^{6-4p-2q} \log(l/144)
\label{eq:men6-4}
\end{align}
(and if needed, here we can use that $\log(l/144) \le \log(l/3)$).
Once more, by induction
\begin{align*}
%\notag
%\left| \manF(2^{m_1}3^{m_2}F_{n_3}^{m_3-1} \cdots F_{n_{k_l}}^{m_{k_l}}) \right|
S_5 
&\le \frac{C}{2^{r_l-1-pm_1-qm_2}} \frac{(r_l-2)!}{m_1!\,m_2!\,(m_3-1)! \cdots m_{k_l}!} \log(l/F_{n_3}) \\
&\qquad= \frac{C}{2^{r_l-pm_1-qm_2}} \frac{(r_l-2)!}{m_1!\,m_2!\,m_3! \cdots m_{k_l}!} \, 2m_3 \log(l/F_{n_3})
%\label{eq:men6-5}
\end{align*}
up to
\begin{align*}
%\notag
%\left| \manF(2^{m_1}3^{m_2}F_{n_3}^{m_3} \cdots F_{n_{k_l}}^{m_{k_l}-1}) \right|
S_6 
&\le \frac{C}{2^{r_l-1-pm_1-qm_2}} \frac{(r_l-2)!}{m_1!\,m_2!\,m_3! \cdots (m_{k_l}-1)!} \log(l/F_{n_{k_l}}) \\
&\qquad= \frac{C}{2^{r_l-pm_1-qm_2}} \frac{(r_l-2)!}{m_1!\,m_2!\,m_3! \cdots m_{k_l}!} 
\, 2m_{k_l} \log(l/F_{n_{k_l}}).
%\label{eq:men6-6}
\end{align*}

The strategy at this point is to combine the coefficients of \eqref{eq:men6-1} and \eqref{eq:men6-2} and to choose $p$ so as to obtain the inequality $\le 2m_1$, then do the same with \eqref{eq:men6-3} and \eqref{eq:men6-4}, choosing $q$ so as to obtain the inequality~$\le 2m_2$.
% La idea ahora es agrupar los coeficientes de \eqref{eq:men6-1} y \eqref{eq:men6-2} 
% y elegir $p$ para que salga un $\le 2m_1$;
% y lo mismo con \eqref{eq:men6-3} y \eqref{eq:men6-4} y elegir $q$ para que salga un $\le 2m_2$.
Specifically, suppose that after using $\log(l/8) \le \log(l/2)$ and $\log(l/144) \le \log(l/3)$, we are able to obtain in \eqref{eq:men6-1} and \eqref{eq:men6-2}, that
% En concreto, supongamos que, tras hacer $\log(l/8) \le \log(l/2)$ y $\log(l/144) \le \log(l/3)$, 
% en \eqref{eq:men6-1} y \eqref{eq:men6-2} conseguimos
\begin{equation}
\label{eq:pm1}
2^{1-p} m_1 + \frac{m_1(m_1-1)(m_1-2)}{(r_l-2)(r_l-3)} \, 2^{3-3p} \le 2m_1
\end{equation}
and likewise, that in \eqref{eq:men6-3} and \eqref{eq:men6-4} we are able to obtain
\begin{equation}
\label{eq:qm2}
2^{1-q} m_2 + \frac{m_1(m_1-1)(m_1-2)(m_1-3)m_2(m_2-1)}{(r_l-2)(r_l-3)(r_l-4)(r_l-5)(r_l-6)} 
\, 2^{6-4p-2q} \le 2m_2.
\end{equation}
If we achieve this, then we would have (possibly with $m_2=0$)
\begin{align*}
|\manF(l)| &\le
\frac{1}{2} \frac{C}{2^{r_l-pm_1-qm_2}} \frac{(r_l-2)!}{m_1!\,m_2!\,m_3! \cdots m_{k_l}!} \\
&\qquad \times
\left( 2m_1\log(l/2) + 2m_2\log(l/3) + 2m_3\log(l/F_{n_3}) + \cdots + 2m_{k_l}\log(l/F_{n_{k_l}}) \right)
\\
&= \frac{1}{2} \frac{C}{2^{r_l-pm_1-qm_2}} \frac{(r_l-2)!}{m_1!\,m_2!\,m_3! \cdots m_{k_l}!}
\log\bigg( \frac{l^{2m_1+2m_2+2m_3+\cdots+2m_{k_l}}}{2^{2m_1}3^{2m_2}F_{n_3}^{2m_3} 
  \cdots F_{n_{k_l}}^{2m_{k_l}}} \bigg)
\\
&= \frac{1}{2} \frac{C}{2^{r_l-pm_1-qm_2}} \frac{(r_l-2)!}{m_1!\,m_2!\,m_3! \cdots m_{k_l}!}
\log(l^{2r_l-2})
\\
&= \frac{C}{2^{r_l-pm_1-qm_2}} \frac{(r_l-1)!}{m_1!\,m_2!\,m_3! \cdots m_{k_l}!}
\log(l),
\end{align*}
and we could conclude.

Let us then deal with \eqref{eq:pm1} and \eqref{eq:qm2}.
In \eqref{eq:pm1}, $r_l \ge m_1+1$, unless $l=2^{m_1}$, in which case we are in the situation of Proposition~\ref{prop:Man(2m)} and there is nothing to prove; on top of this we already said that we need to take $p=1$ even if we discard that case, hence
\begin{equation}
\label{eq:m1rl}
\frac{(m_1-1)(m_1-2)}{(r_l-2)(r_l-3)} \le 1,
\end{equation}
and what we need is to find $p \in [0,1]$ such that $2^{1-p} + 2^{3-3p} \le 2$, or equivalently,
\begin{equation}
\label{eq:f(p)}
f(p) := 2^{-p} + 2^{2-3p} \le 1.
\end{equation}
The function $f(p)$ is decreasing on $[0,1]$, with $f(0)=5$ and $f(1)=1$, which leaves us only with the choice $p=1$. We could choose a value $p < 1$ if $m_{1}$ and $r_{1}$ were restricted in such a way as to yield a smaller bound in \eqref{eq:m1rl}, although then we would not have the same function $f(p)$ as above in \eqref{eq:f(p)} (Remark~\ref{rem:m1}, goes into more detail regarding this possibility).

In general we cannot do this, since if $m_{1}$ and $r_{l}$ both go to $\infty$ in such a way that their quotient goes to $1$, it is impossible to obtain a smaller bound in \eqref{eq:m1rl}. However, it would be possible if we assume a relation along the lines of $r_l \ge 2m_1$ (see again Remark~\ref{rem:m1}).

In any case, with $p=1$, after canceling $2m_2$, \eqref{eq:qm2} can be rewritten as
% (fijarse que he subido un $2$ al numerador)
\begin{equation}
\label{eq:qm2-sinp}
2^{-q} + \frac{2m_1(m_1-1)(m_1-2)(m_1-3)(m_2-1)}{(r_l-2)(r_l-3)(r_l-4)(r_l-5)(r_l-6)} \, 2^{-2q} \le 1.
\end{equation}
Since this bound involves $\IF_{144}$, we have $m_1\ge 4$ y $m_2\ge 2$, so that we may assume that $r_l \ge 7$ (if $r_l = 6$ then $l=188$ and one has to verify separately that this case also satisfies the bound).
%% En principio, podr\'ia ocurrir que satisface con constante multiplicativa mayor si hiciera falta, 
%% pero tambi\'en se puede pasar completamente de este caso pues, para estudiar la convergencia de una serie, 
%% un sumando no pinta nada.
Now, with $m_1\ge 4$, $m_2\ge 2$, $r_l \ge 7$, it is easy to check that, except when $r_l=m_1+m_2$, with
\begin{equation}
\label{eq:excepciones}
(m_1,m_2) \in \big\{ (5,2), (6,2), (7,2), (4,3), (5,3), (6,3) \big\},
\end{equation}
we have
\begin{equation}
\label{eq:1/2}
\frac{2m_1(m_1-1)(m_1-2)(m_1-3)(m_2-1)}{(r_l-2)(r_l-3)(r_l-4)(r_l-5)(r_l-6)} \le \frac{1}{2}.
\end{equation}

Since the exceptions in \eqref{eq:excepciones} are a finite number of cases, we can simply consult Table~\ref{tab:m1m2} to check that in each of these cases the bound is satisfied with $C=7/4$ as we stated at the beginning.
%% No tienen ninguna importancia para decidir si la serie converge o no, 
%% o porque s\'olo afectar\'ia al valor de~$C$
Thus, what we need to do now is to find $q \in [0,1]$ such that
\begin{equation}
\label{eq:g(q)}
g(q) := 2^{-q} + 2^{-2q-1} \le 1.
\end{equation}
This is a decreasing function on $[0,1]$, with $g(0)=3/2$ and $g(1)=5/8$; to obtain
$g(q)\le1$ it is enough to take
\[
q = \log_2\bigg(\frac{1+\sqrt{3}}{2}\Big) = \log\Big(\frac{1+\sqrt{3}}{2}\bigg)/\log(2)
= 0.449984\dots
\]
or anything larger. 

Without going into all the details, let us see how this method can be improved to obtain a smaller $q$. If in \eqref{eq:1/2} we substitute $\le 1/2$ for $\le 2/5$, a finite number of exceptions appear: we need to add $(7,3)$, $(8,2)$ y $(8,3)$ to those in \eqref{eq:excepciones}); the inequality which replaces \eqref{eq:g(q)} is now $g(q) := 2^{-q} + 2^{-2q-1}/5 \le 1$, and we can lower $q$ down to $q = \log_2\Big(\frac{5+\sqrt{65}}{10}\Big) = 0.3854\dots$.
%% En todo caso, esto es muy f\'acil de chequear, pues el Reduce de Mathematica da todas las excepciones.

With $\le 1/4$ in \eqref{eq:1/2}, somewhat more than 30 exceptions arise; this time, there are 3 exceptions with $r_l=m_1+m_2+1$, which are easily dealt with with a computational algebra package. Thus leads to $q = \log_2\Big(\frac{1+\sqrt{2}}{2}\Big) = 0.271553\dots$

With $\le1/5$ now there are a bit more than 100 exceptions, and we continue lowering the value of $q$. This time we have the inequality $g(q) := 2^{-q} + 2^{-2q}/5 \le 1$, and we arrive at
\[
  q = \log_2 \bigg( \frac{1}{2} + \frac{3}{10} \sqrt{5} \bigg) = 0.22752\dots,
\]
which concludes the proof of the theorem. Note that with $\le1/6$, an infinite number of exceptions arise, so we cannot continue.
\end{proof}

\begin{remark}
\label{rem:m1}
If we look for bounds under the condition that $r_{l} \geq 2m_{1}$, then instead of \eqref{eq:m1rl} we have
\begin{equation*}
%\label{eq:m1rl4}
\frac{(m_1-1)(m_1-2)}{(r_l-2)(r_l-3)} \le \frac{1}{4},
\end{equation*}
which leads to finding $p \in [0,1]$ such that $2^{1-p} + 2^{1-3p} \le 2$, i.e.
\begin{equation*}
%\label{eq:f(p)alt}
f(p) := 2^{-p} + 2^{-3p} \le 1.
\end{equation*}
The unique solution $f(p)=1$ with $p \in [0,1]$ is
\[
p = -\log_2\bigg( \sqrt[3]{ \frac{1}{18} (\sqrt{93}+9) }
- \sqrt[3]{ \frac{1}{18} (\sqrt{93}-9)} \bigg)
= 0.551463\dots
\]
Note however that this does not just consist of a slight variation on the proof, since the condition $r_l \ge 2m_1$ doesn't play well with our inductive steps. Other ideas  are needed to get it to work. In addition, this would also affect $q$ in \eqref{eq:qm2} and would not lead to the same expression as in~\eqref{eq:qm2-sinp}.
\end{remark}

%----------
\section{Proof of Theorem~\ref{theo:sigma0}}
%----------

We prove that the Dirichlet series \eqref{eq:ManDir} does not converge when $s$ satisfies~\eqref{eq:diverge}, by showing that the general term $\manF(l) l^{-s}$ does not tend to~$0$. Without loss of generality, we can consider that $s$ is a positive real number.

The idea is to take $l$ a product of consecutive Fibonacci numbers greater than $2$ (excluding $F_6$ and $F_{12}$ of course) to a power $m$, get the value $\manF(l)$ applying Proposition~\ref{prop:Man(l)prod}, estimate it by using Stirling's formula, and obtain the condition for $\manF(l)l^{-s}$, as a function of $m$, to tend to $0$ when $m\to\infty$. We have checked that the most demanding condition on $s$ is obtained with seven Fibonacci factors, so we will limit ourselves to this case.

Let $l = (3\cdot 5\cdot 13\cdot 21\cdot 34\cdot 55\cdot 89)^m$, with $m \ge 1$. By Proposition~\ref{prop:Man(l)prod},
\[
\manF(l) = (-1)^{7m-1} \frac{1}{2^{7m}} \frac{(7m-1)!}{(m!)^7}\, m 
\log(3\cdot 5\cdot 13\cdot 21\cdot 34\cdot 55\cdot 89).
\]
Stirling’s formula leads to
%% With Mathematica:
%% Asymptotic[(7*m-1)!, m -> Infinity]
%% y dividir "a mano" por
%% Asymptotic[m!, m -> Infinity]
%% elevado a la 7
%% (la asintótica de Asymptotic[(7*m-1)!/(m!)^7, m -> Infinity] la devuelve sin hacer)
\[
\manF(l) l^{-s} \sim (-1)^{7m-1} \left(\frac{7}{2}\right)^{7m} \frac{1}{8\sqrt{7}\pi^3 m^3} 
\log(3\cdot 5\cdot 13\cdot 21\cdot 34\cdot 55\cdot 89) \cdot
(3\cdot 5\cdot 13\cdot 21\cdot 34\cdot 55\cdot 89)^{-sm},
\quad m \to \infty.
\]
Convergence of the series $\sum_{l} \manF(l) l^{-s}$ implies that
\[
\frac{1}{m^3} \left(\frac{7}{2}\right)^{7m} 
(3\cdot 5\cdot 13\cdot 21\cdot 34\cdot 55\cdot 89)^{-sm} \to 0 \quad \text{when }m \to \infty
\]
and it is an easy exercise to check that this is not true when
\[
s < \frac{7\log(7/2)}{\log(3\cdot 5\cdot 13\cdot 21\cdot 34\cdot 55\cdot 89)} = 0.431141\dots,
\]
so that \eqref{eq:ManDir} does not converge in this case.
%% Sum[1/2^(7m) (7m-1)!/(m!)^7*m*(3*5*13*21*34*55*89)^{-s*m}, {m, 1, Infinity}, GenerateConditions -> True]
%% La condición es 681530850^s > 823543/128, i.e. 
%% s > \frac{7\log(7/2)}{\log(3\cdot 5\cdot 13\cdot 21\cdot 34\cdot 55\cdot 89)}
%% Reduce[681530850^s > 823543/128, s, Reals]
%% s > -((7 (Log[2] - Log[7]))/(Log[2] + 2 Log[3] + 2 Log[5] + Log[7] + Log[11] + Log[13] + Log[17] + Log[89]))
%% Con el = el t\'ermino general s\'i se va a 0, pero la serie que le he planteado a Mathematica no converge.

We will now find an abscissa $\sigma_0$ such that $\sum_l {\manF(l)} l^{-s}$ converges absolutely for $\Re(s)>\sigma_0$, necessarily to $\varphi'(s)/\varphi(s)$. To achieve this, we use the bound in Proposition~\ref{prop:Mcg}, which will show how the resulting value of $\sigma_{0}$ depends on the numbers $p,q$ which appear in~\eqref{eq:Mcg}. The smaller we can take $p$ and $q$, the smaller the value of~$\sigma_0$. 

Considering $\sum_l \frac{\manF(l)}{\log(l)} l^{-s}$ instead of $\sum_l {\manF(l)} l^{-s}$ does not change the abscissa of absolute convergence (it can change what happens on the boundary, but this doesn't affect our argument), hence we can drop the factors of $\log(l)$. In addition, since our bounds are absolute, we are actually dealing with  $\sum_l {|\manF(l)|} l^{-s}$ for positive real $s$.

With this in mind, consider the function
\[
\xi(s) = \sum_{l \in \Fp} \frac{1}{2^{r_l-pm_1-qm_2}} \frac{1}{r_l} 
\frac{r_l!}{m_1!\,m_2!\,m_3! \cdots m_{k_l}!} \,l^{-s}
\]
where $l \in \Fp$ has the factorization $l = F_{n_1}^{m_1}F_{n_2}^{m_2}F_{n_3}^{m_3}  \cdots F_{n_{k_l}}^{m_{k_l}}$ given in \eqref{eq:Fp}. Since
\[
\sum_{l \in \Fp} |\manF(l)/\log(l)| l^{-s} \le \xi(s),
\quad 
\]
the convergence of $\xi(s)$ implies the absolute convergence of $\sum_l {\manF(l)} l^{-s}$ to $\varphi'(s)/\varphi(s)$. 

For a positive integer $M$, consider the set of $l \in \Fp$ whose factorization~\eqref{eq:Fp} contains only Fibonacci numbers from $F_{3} = 2$ up to $F_{M}$ (some may be absent in the factorization), in other words,
\begin{align*}
%\label{eq:FpM}
\Fp(M) := \big\{\,
l &= F_{n_1}^{m_1}F_{n_2}^{m_2}F_{n_3}^{m_3} \cdots F_{M}^{m_{M}} : m_{j} \ge 0, \\ 
&\qquad\qquad 3 = n_1<n_2<\cdots < n_{k_l} = M,  
n_{j}\neq 6, 12,\, n_j \le M \,\big\}.
\end{align*}
Assuming $M>12$, we have
\begin{equation*}
%\label{eq:nj}
\begin{gathered}
n_1 = 3,\
n_2 = 4,\
n_3 = 5,\
n_4 = 7,\
n_5 = 8,\
n_6 = 9,\
n_7 = 10,\
n_8 = 11,\
n_9 = 13,\\
n_j = j+4\text{ for }j > 9,
\quad\text{up to}\quad
n_{k_l} = M-4.
\end{gathered}
\end{equation*}
Consider the summands of the series $\xi(s)$ corresponding to $l \in \Fp(M)$,
\begin{align*}
\xi_M(s) &= \sum_{l \in \Fp(M)} \frac{1}{2^{r_l-pm_1-qm_2}} 
  \frac{1}{r_l} \frac{r_l!}{m_1!\,m_2!\,m_3! \cdots m_{k_l}!} 
  \Big(F_{n_1}^{m_1}F_{n_2}^{m_2}F_{n_3}^{m_3} \cdots F_{n_{k_l}}^{m_{k_l}}\Big)^{-s} \\
  &= \sum_{l \in \Fp(M)} \frac{1}{r_l} \frac{r_l!}{m_1!\,m_2!\,m_3! \cdots m_{k_l}!} 
  \Big((2^{1-p}F_{n_1}^{-s})^{m_1} (2^{1-q}F_{n_2}^{-s})^{m_2} (2F_{n_3}^{-s})^{m_3} \cdots (2F_{n_{k_l}}^{-s})^{m_{k_l}}\Big).
\end{align*} 
We then have $\xi(s) = \lim_{M\to\infty} \xi_M(s)$.
% Obs\'ervese que, en las factorizaciones de $l \in \Fp(M)$, podemos incluir todos los $F_{n_j} \le M$ 
% (con $n_{j}\neq 6, 12$) sin m\'as que poner $m_j \ge 0$ en vez de $m_j > 0$; 
% los $m_j = 0$ no tienen ninguna contribuci\'on ni en $r_l$ ni en el producto de los $m_j!$.
For fixed $r$, let
\[
S_{M,r}(s) = \frac{1}{r} \sum_{\substack{l \in \Fp(M)\\ r_l=r}}  
\frac{r_l!}{m_1!\,m_2!\,m_3! \cdots m_{k_l}!} 
  \Big((2^{1-p}F_{n_1}^{-s})^{m_1} (2^{1-q}F_{n_2}^{-s})^{m_2} 
  (2F_{n_3}^{-s})^{m_3} \cdots (2F_{n_{k_l}}^{-s})^{m_{k_l}}\Big).
\]
% Since indices with $m_j = 0$ do not contribute to the product of the $m_{j}!$ or to $r_{l}$, 
% the formula for $\manF(l)$ sigue siendo v\'alida.
The multinomial expansion
\[
(x_1+x_2+\cdots+x_k)^{r} =
\sum_{\substack{m_1+m_2+\cdots+m_k=r,\\ m_j \ge 0}} \frac{r!}{m_1!\,m_2!\cdots m_k!} \, x_1^{m_1} x_2^{m_2} \cdots x_k^{m_k}
\]
leads directly to
\[
S_{M,r}(s) = \frac{1}{r} \left( \frac{1}{2^{1-p}F_3^s} + \frac{1}{2^{1-q}F_4^s} + \frac{1}{2F_5^s} 
+ \dots + \frac{1}{2F_{n_{k_l}}^s} \right)^r
=: \frac{1}{r} \Big( S_{M}^{*}(s) \Big)^r,
\]
where the $F_{j}$ appearing in this expression are all those $F_j \le F_M$ with $j \notin \{1,2,6,12\}$ 

Since all the terms involved are positive for $s > 0$, Tannery's Theorem, which is a special case of the Dominated Convergence Theorem, shows that
\[
  \xi(s) = \lim_{M\to\infty} \xi_M(s)
  = \lim_{M\to\infty} \sum_{r=1}^{\infty} S_{M,r}(s)
  = \sum_{r=1}^{\infty} \lim_{M\to\infty} S_{M,r}(s)
  = \sum_{r=1}^{\infty} \frac{1}{r} \Big(\lim_{M\to\infty} S_{M}^{*}(s) \Big)^r,
\]
where the limit is exchanged with raising to the $r$\textsuperscript{th} power is justified by continuity. Finally,
\begin{align*}
&\lim_{M\to\infty} S_{M}^{*}(s)
= \frac{1}{2^{1-p}F_3^s} + \frac{1}{2^{1-q}F_4^s} + \frac{1}{2F_5^s} + \cdots + \frac{1}{2F_{M}^s} + \cdots \\
&= \frac{1}{2^{1-p}F_3^s} + \frac{1}{2^{1-q}F_4^s} + \frac{1}{2} \sum_{j=5}^{\infty} \frac{1}{F_j^s}
- \frac{1}{2F_6^s} - \frac{1}{2F_{12}^s} \\
&= \frac{1}{2^{1-p}2^s} + \frac{1}{2^{1-q}3^s} + \frac{1}{2} \varphi(s) - 1
- \frac{1}{2 \cdot 2^s} - \frac{1}{2 \cdot 3^s} - \frac{1}{2 \cdot 8^s} - \frac{1}{2 \cdot 144^s} 
=: \varphi^{*}(s).
\end{align*}
Recall that $\sum_{r=1}^\infty a^r/r = -\log(1-a)$ for $|a|<1$. Hence, if $\varphi^{*}(s) < 1$, then
\[
\xi(s) = \sum_{r=1}^{\infty} \frac{1}{r} \Big(\varphi^{*}(s)\Big)^r
= -\log(1-\varphi^{*}(s))
\]
and we will be done. Since $\varphi^{*}(s)$ is clearly decreasing, $\varphi^{*}(s) < 1$ will hold for $s > \sigma_0$ for that $\sigma_0$ satisfying $\varphi^{*}(\sigma_0) = 1$; the concrete value depends on $p$ and $q$. With $p=1$ and $q = 0.22752\dots$ as in Proposition~\ref{prop:Mcg}, we find $\sigma_0=0.905556\dots$, which concludes the proof.

\begin{remark}
As we saw in the proof of Proposition~\ref{prop:Mcg}, the bound in \eqref{eq:Mcg} with $p=q=0$ is false, but, if it were valid, the corresponding solution of $\varphi^{*}(s) = 1$ would be $s=0.686895\dots$. The bound above is probably not optimal, and we can speculate that the abscisa of absolute convergence of \eqref{eq:ManFp} is some number between our $\sigma_0=0.905556\dots$, obtained in the proof of Theorem~\ref{theo:sigma0}, and the value $0.686895\dots$ which would result if the bound \eqref{eq:Mcg} with $p=q=0$ held.
\end{remark}

%----------
\begin{thebibliography}{99}
%----------

\bibitem{Ap} 
T. M. Apostol,  
\textit{Introduction to Analytic Number Theory}, 
Springer, New York, 1976.

%\bibitem{Ash} 
%Ash, R. B.: 
%\textit{Complex Variables}, 
%Academic Press, New York, 1971.

%\bibitem{Ber}
%V. Bernstein, % Vladimir Bernstein
%\textit{Le\c{c}ons sur les progr\`es r\'ecents de la th\'eorie des s\'eries 
%de Dirichlet profess\'ees au Coll\`ege de France},
%Gauthier-Villars, Paris, 1933.
%% 320 pages
%% https://books.google.es/books?id=UnB2oAEACAAJ
%% https://zbmath.org/0008.11503

% MR2399185 
\bibitem{BuLuMiSi} 
% Bugeaud, Y., Luca, F., Mignotte, M., Siksek, S.: 
Y. Bugeaud, F. Luca, M. Mignotte and S. Siksek, 
{Fibonacci numbers at most one away from a perfect power}, 
\textit{Elem. Math.} \textbf{63} (2008), 65--75.

% MR2215137
\bibitem{BMS}
Y. Bugeaud, M. Mignotte and S. Siksek,
Classical and modular approaches to exponential Diophantine equations I. 
Fibonacci and Lucas perfect powers, 
\textit{Ann. of Math. (2)} \textbf{163} (2006), 969--1018.

% MR1502458 y MR1502459 
\bibitem{Car}
R. D. Carmichael, 
On the numerical factors of the arithmetic forms $\alpha^n\pm \beta^n$,
\textit{Ann. of Math. (2)} \textbf{15} (1913/14), 30--48 and 49--80.

%%% NO HACE FALTA
%% Corolario de p. 262:
%% Si la hip\'otesis de Riemann es verdad, entonces la serie de Dirichet de $1/\zeta(s)$
%% converge a $1/\zeta(s)$ en todo el semiplano en $\Re(s) > 1/2$
%% (El rec\'iproco es verdad y es f\'acil).
%% De hecho, la serie de Dirichlet de $1/\zeta(s)$ es la de coeficientes $\mu(n), 
%% y precisamente la mu_F(n) no sale en este artículo.
% \bibitem{Ed}
% H. M. Edwards,
% \textit{Riemann's Zeta Function},
% Academic Press, 1974.

% MR1744904
\bibitem{Ega} 
S. Egami, % Egami, Shigeki
{Some curious Dirichlet series} (English Summary), 
\textit{Analytic number theory and its interactions with other
parts of number theory} (Japanese) (Kyoto, 1998),
\textit{S\={u}rikaisekikenky\={u}sho K\={o}ky\={u}roku} \textbf{1091} (1999), 172--174.

% MR3589830
\bibitem{GoSh}
H. W. Gould and T. Shonhiwa,
A catalog of interesting Dirichlet series,
\textit{Missouri J. Math. Sci.} \textbf{20} (2008), 1--17. % no. 1 
% https://doi.org/10.35834/mjms/1316032830

\bibitem{HaWr}
G. H. Hardy and E. M. Wright,
\textit{An introduction to the theory of numbers},
5th~ed., Oxford University Press, 1979.

% MR0792089 y MR1994094
\bibitem{Ivic} 
% Ivi\'c, Aleksandar
A. Ivi\'c,  
\textit{The Riemann zeta-function}, John Wiley \& Sons,
New York, 1985. Reprinted, Dover, 2003.

%%% Son dos p\'aginas y es majo, 
%%% pero no est\'a publicado y tampoco aporta nada importante.
%\bibitem{Kit} 
%A. R. Kitson, % Adrian R. Kitson
%{The regularized product of the Fibonacci numbers}, 
%arXiv:math/0608187, 8 Aug 2006.
%% https://arxiv.org/abs/math/0608187

\bibitem{Kno}
J. Knopfmacher, % John Knopfmacher
\textit{Abstract Analytic Number Theory}, Dover, 2015.
% Corrected and enlarged republication of the edition originally published 
% by North-Holland Publishing Company, Amsterdam, 
% and the American Elsevier Publishing Company, Inc., New York, 1975.

\bibitem{Mo}
G. Mora,
On the zeros of the partial sums of the Fibonacci zeta function,
\textit{Rev. R. Acad. Cienc. Exactas F\'is. Nat. Ser. A Mat. RACSAM} 
\textbf{117} (2023), no.~4, Paper No. 139, 18~pp.
% https://doi.org/10.1007/s13398-023-01471-1

% MR3156859
\bibitem{Mur} 
% Murty, M. Ram
M. R. Murty, 
{The Fibonacci zeta function}, 
\textit{Automorphic representations and $L$-functions} (Mumbai, 2012), 409--425, 
Tata Inst. Fundam. Res. Stud. Math., {22}, Hindustan Book Agency, Gurgaon, 2013.

% MR2279934
\bibitem{MurSin} 
% Murty, M. Ram; Sinha, Kaneenika
M. R. Murty and K. Sinha,  
{Multiple Hurwitz zeta functions}, 
\textit{Multiple Dirichlet Series, Automorphic Forms and Analytic Number Theory}, 
135--156, Proc. Sympos. Pure Math., {75}, Amer. Math. Society, Providence, RI, 2006.

\bibitem{Na}
L. M. Navas, 
Analytic continuation of the Fibonacci Dirichlet series, 
\textit{Fibonacci Q.} \textbf{39} (2001), 409--418.
% https://www.fq.math.ca/Scanned/39-5/

%\bibitem{Rob}
%H. Robbins, 
%A remark on Stirling's formula, 
%\textit{Amer. Math. Monthly} \textbf{62} (1955), no.~1, 26--29.
%% https://www.jstor.org/stable/2308012
%% https://doi.org/10.2307%2F2308012

% MR4594520 
\bibitem{SeNa}
\'A. Serrano Holgado and L. M. Navas,
The zeta function of a recurrence sequence of arbitrary degree,
\textit{Mediterr. J. Math.} \textbf{20} (2023), no.~4, Paper No.~224, 18~pp.
% https://doi.org/10.1007/s00009-023-02427-w

\bibitem{Ya}
% Yabuta, Minoru
M. Yabuta, 
A simple proof of Carmichael's theorem on primitive divisors,
\textit{Fibonacci Q.} \textbf{39} (2001), no.~5, 439--443.
% https://www.fq.math.ca/Scanned/39-5/

%----------
\end{thebibliography}
%----------

%----------
\end{document}
%----------

